\documentclass[10pt,a4paper]{amsart}
\usepackage{amsmath,amssymb,amsthm,graphicx}
% For commutative diagrams we use Xy-pic
\usepackage[all]{xy}

% We use 2cell for 2-commutative diagrams.
\xyoption{2cell}
% Theorem environments.
%
\theoremstyle{plain}
\newtheorem{theorem}[subsection]{Theorem}
\newtheorem{proposition}[subsection]{Proposition}
\newtheorem{lemma}[subsection]{Lemma}

\theoremstyle{definition}
\newtheorem{definition}[subsection]{Definition}
\newtheorem{example}[subsection]{Example}
\newtheorem{exercise}[subsection]{Exercise}
\newtheorem{situation}[subsection]{Situation}

\theoremstyle{remark}
\newtheorem{remark}[subsection]{Remark}
\newtheorem{remarks}[subsection]{Remarks}

\numberwithin{equation}{subsection}

% Macros
%
\def\lim{\mathop{\mathrm{lim}}\nolimits}
\def\colim{\mathop{\mathrm{colim}}\nolimits}
\def\Spec{\mathop{\mathrm{Spec}}}
\def\Hom{\mathop{\mathrm{Hom}}\nolimits}
\def\Ext{\mathop{\mathrm{Ext}}\nolimits}
\def\SheafHom{\mathop{\mathcal{H}\!\mathit{om}}\nolimits}
\def\SheafExt{\mathop{\mathcal{E}\!\mathit{xt}}\nolimits}
\def\Sch{\mathit{Sch}}
\def\Mor{\mathop{\mathrm{Mor}}\nolimits}
\def\Ob{\mathop{\mathrm{Ob}}\nolimits}
\def\Sh{\mathop{\mathit{Sh}}\nolimits}
\def\NL{\mathop{N\!L}\nolimits}
\def\CH{\mathop{\mathrm{CH}}\nolimits}
\def\proetale{{pro\text{-}\acute{e}tale}}
\def\etale{{\acute{e}tale}}
\def\QCoh{\mathit{QCoh}}
\def\Ker{\mathop{\mathrm{Ker}}}
\def\Im{\mathop{\mathrm{Im}}}
\def\Coker{\mathop{\mathrm{Coker}}}
\def\Coim{\mathop{\mathrm{Coim}}}

% Boxtimes
%
\DeclareMathSymbol{\boxtimes}{\mathbin}{AMSa}{"02}

%
% Macros for moduli stacks/spaces
%
\def\QCohstack{\mathcal{QC}\!\mathit{oh}}
\def\Cohstack{\mathcal{C}\!\mathit{oh}}
\def\Spacesstack{\mathcal{S}\!\mathit{paces}}
\def\Quotfunctor{\mathrm{Quot}}
\def\Hilbfunctor{\mathrm{Hilb}}
\def\Curvesstack{\mathcal{C}\!\mathit{urves}}
\def\Polarizedstack{\mathcal{P}\!\mathit{olarized}}
\def\Complexesstack{\mathcal{C}\!\mathit{omplexes}}
% \Pic is the operator that assigns to X its picard group, usage \Pic(X)
% \Picardstack_{X/B} denotes the Picard stack of X over B
% \Picardfunctor_{X/B} denotes the Picard functor of X over B
\def\Pic{\mathop{\mathrm{Pic}}\nolimits}
\def\Picardstack{\mathcal{P}\!\mathit{ic}}
\def\Picardfunctor{\mathrm{Pic}}
\def\Deformationcategory{\mathcal{D}\!\mathit{ef}}
\usepackage{cleveref}
\crefname{originallemma}{original lemma}{original lemmas}
\Crefname{originallemma}{Original lemma}{Original lemmas}
\makeatletter
\newlabel{lemma-base-change-q-injective}{{0F03}{source}{}{sourcepointer-1}{}}
\newlabel{lemma-base-change-q-injective@cref}{{[originallemma][0][]0F03}{[1][0][]source}{}{}{}}
\makeatother
\crefname{originallemma}{original lemma}{original lemmas}
\Crefname{originallemma}{Original lemma}{Original lemmas}
\makeatletter
\newlabel{lemma-formal-argument}{{0F06}{source}{}{sourcepointer-2}{}}
\newlabel{lemma-formal-argument@cref}{{[originallemma][0][]0F06}{[1][0][]source}{}{}{}}
\makeatother
\crefname{originalproposition}{original proposition}{original propositions}
\Crefname{originalproposition}{Original proposition}{Original propositions}
\makeatletter
\newlabel{proposition-topological-invariance}{{03SI}{source}{}{sourcepointer-3}{}}
\newlabel{proposition-topological-invariance@cref}{{[originalproposition][0][]03SI}{[1][0][]source}{}{}{}}
\makeatother
\crefname{originallemma}{original lemma}{original lemmas}
\Crefname{originallemma}{Original lemma}{Original lemmas}
\makeatletter
\newlabel{lemma-relative-colimit-general}{{0EYM}{source}{}{sourcepointer-4}{}}
\newlabel{lemma-relative-colimit-general@cref}{{[originallemma][0][]0EYM}{[1][0][]source}{}{}{}}
\makeatother
\crefname{originalproposition}{original proposition}{original propositions}
\Crefname{originalproposition}{Original proposition}{Original propositions}
\makeatletter
\newlabel{proposition-finite-higher-direct-image-zero}{{03QP}{source}{}{sourcepointer-5}{}}
\newlabel{proposition-finite-higher-direct-image-zero@cref}{{[originalproposition][0][]03QP}{[1][0][]source}{}{}{}}
\makeatother
\crefname{originallemma}{original lemma}{original lemmas}
\Crefname{originallemma}{Original lemma}{Original lemmas}
\makeatletter
\newlabel{lemma-finite-pushforward-commutes-with-base-change}{{0959}{source}{}{sourcepointer-6}{}}
\newlabel{lemma-finite-pushforward-commutes-with-base-change@cref}{{[originallemma][0][]0959}{[1][0][]source}{}{}{}}
\makeatother
\crefname{originallemma}{original lemma}{original lemmas}
\Crefname{originallemma}{Original lemma}{Original lemmas}
\makeatletter
\newlabel{sites-cohomology-lemma-relative-Leray}{{0734}{source}{}{sourcepointer-7}{}}
\newlabel{sites-cohomology-lemma-relative-Leray@cref}{{[originallemma][0][]0734}{[1][0][]source}{}{}{}}
\makeatother
\title{Failure of base change in degree q produces a field-valued counterexample}
\author{The Stacks Project Authors}
\date{stacks-project:etale-cohomology.tex:local-snapshot-20261011; GNU FDL 1.2 or later}
\begin{document}
\maketitle
\noindent\textit{Editorial note.} All delivered mathematical text is the author's own, copied from the pinned Stacks Project chapter file etale-cohomology.tex; no statement or proof step was rewritten, reordered or removed.\par\medskip
\noindent\textit{Editorial note.} The counted statement begins with 'With f : X -> S and n as in Remark ...', so the remark introducing BC(f, n, q\_0) is delivered as a context span (source lines 16847-16866); it is the author's own text copied verbatim and makes the card self-contained.\par\medskip
\noindent\textit{Editorial note.} The proof is the author's complete verification of the three hypotheses, including the decomposition step which the author singles out as the most interesting; the final appeal to the formal argument is the author's own and the formal argument lemma itself is cited as a pointer rather than reproduced.\par\medskip
\noindent\textit{Editorial note.} Every labelled reference occurring inside the counted spans points to a chapter of the Project that is not part of this card's source tree. Each is carried in assembly.external\_references as a non-load-bearing navigation pointer whose source location is the invocation inside the counted span, so no load-bearing mathematics is replaced by a stub.\par\medskip
\noindent\textit{Editorial note.} The preamble is the chapter's own native preamble excerpt already used for this source: the Xy-pic setup (native preamble lines 14-18) and the native theorem declarations with the mathematical macros (native preamble lines 154-217). Both Xy-pic diagrams are delivered as written.\par\medskip
\noindent\textit{Allowed standard prerequisites (editorial).}\begin{itemize}
\item The property BC(f, n, q\_0) of a morphism of schemes. Same chapter, the remark defining BC(f, n, q\_0): for every commutative square with g quasi-compact and quasi-separated and every abelian sheaf annihilated by n on T\_etale, the base change map is an isomorphism in degrees at most q\_0. Supplies the notation of the statement and the property whose failure is upgraded to a field-valued counterexample.
\item Invariance of the etale topology under thickenings. Same chapter, the topological invariance of the etale site under a nilpotent thickening. Supplies condition (1) of the formal argument lemma, namely that the property under consideration is thickening-invariant.
\item Cohomology commutes with limits of quasi-compact quasi-separated schemes. Same chapter, the results on limits of schemes and the behaviour of base change under limits, used to establish condition (2) of the formal argument lemma. Supplies the limit-stability of the property P defined in the proof.
\item Descent of the base change property along a closed subscheme with quasi-compact complement. Same chapter, the decomposition step in which a counterexample on T' is restricted to a closed subscheme Z or its complement V. Supplies condition (3) of the formal argument lemma, which the author calls the most interesting one.
\end{itemize}\medskip
\ifdefined\hypertarget\hypertarget{sourcepointer-1}{}\fi
\noindent\textit{Source reference.} Original lemma 0F03: sources/stacks-etale-cohomology/pinned-original-source.tex, source lines 16869--16869. Navigation pointer invoked at proof line 17111; the counted span contains the pointer only and the cited result is not reproduced inside it, so the pointer is registered as a non-load-bearing navigation stub in the established form for Stacks material. No bibliography citation is involved.\par\medskip
\ifdefined\hypertarget\hypertarget{sourcepointer-2}{}\fi
\noindent\textit{Source reference.} Original lemma 0F06: sources/stacks-etale-cohomology/pinned-original-source.tex, source lines 17046--17046. Navigation pointer invoked at proof line 17157; the counted span contains the pointer only and the cited result is not reproduced inside it, so the pointer is registered as a non-load-bearing navigation stub in the established form for Stacks material. No bibliography citation is involved.\par\medskip
\ifdefined\hypertarget\hypertarget{sourcepointer-3}{}\fi
\noindent\textit{Source reference.} Original proposition 03SI: sources/stacks-etale-cohomology/pinned-original-source.tex, source lines 6548--6548. Navigation pointer invoked at proof line 17164; the counted span contains the pointer only and the cited result is not reproduced inside it, so the pointer is registered as a non-load-bearing navigation stub in the established form for Stacks material. No bibliography citation is involved.\par\medskip
\ifdefined\hypertarget\hypertarget{sourcepointer-4}{}\fi
\noindent\textit{Source reference.} Original lemma 0EYM: sources/stacks-etale-cohomology/pinned-original-source.tex, source lines 7567--7567. Navigation pointer invoked at proof line 17192; the counted span contains the pointer only and the cited result is not reproduced inside it, so the pointer is registered as a non-load-bearing navigation stub in the established form for Stacks material. No bibliography citation is involved.\par\medskip
\ifdefined\hypertarget\hypertarget{sourcepointer-5}{}\fi
\noindent\textit{Source reference.} Original proposition 03QP: sources/stacks-etale-cohomology/pinned-original-source.tex, source lines 8245--8245. Navigation pointer invoked at proof line 17252; the counted span contains the pointer only and the cited result is not reproduced inside it, so the pointer is registered as a non-load-bearing navigation stub in the established form for Stacks material. No bibliography citation is involved.\par\medskip
\ifdefined\hypertarget\hypertarget{sourcepointer-6}{}\fi
\noindent\textit{Source reference.} Original lemma 0959: sources/stacks-etale-cohomology/pinned-original-source.tex, source lines 8292--8292. Navigation pointer invoked at proof line 17259; the counted span contains the pointer only and the cited result is not reproduced inside it, so the pointer is registered as a non-load-bearing navigation stub in the established form for Stacks material. No bibliography citation is involved.\par\medskip
\ifdefined\hypertarget\hypertarget{sourcepointer-7}{}\fi
\noindent\textit{Source reference.} Original lemma 0734: sources/stacks-etale-cohomology/pinned-original-source.tex, source lines 17280--17280. Navigation pointer invoked at proof line 17280; the counted span contains the pointer only and the cited result is not reproduced inside it, so the pointer is registered as a non-load-bearing navigation stub in the established form for Stacks material. No bibliography citation is involved.\par\medskip
\begin{remark}
\label{remark-base-change-holds}
Let $f : X \to S$ be a morphism of schemes. Let $n$ be an integer.
We will say $BC(f, n, q_0)$ is true if for every commutative diagram
$$
\xymatrix{
X \ar[d]_f & X' \ar[l] \ar[d]_{f'} & Y \ar[l]^h \ar[d]^e \\
S & S' \ar[l] & T \ar[l]_g
}
$$
with $X' = X \times_S S'$ and $Y = X' \times_{S'} T$ and
$g$ quasi-compact and quasi-separated, and every abelian sheaf
$\mathcal{F}$ on $T_\etale$ annihilated by $n$ the base change map
$$
(f')^{-1}R^qg_*\mathcal{F}
\longrightarrow
R^qh_*e^{-1}\mathcal{F}
$$
is an isomorphism for $q \leq q_0$.
\end{remark}

\begin{lemma}
\label{lemma-base-change-does-not-hold-pre}
With $f : X \to S$ and $n$ as in Remark \ref{remark-base-change-holds}
assume for some $q \geq 1$ we have that $BC(f, n, q - 1)$ is true, but
$BC(f, n, q)$ is not. Then there exist a commutative diagram
$$
\xymatrix{
X \ar[d]_f & X' \ar[d]_{f'} \ar[l] & Y \ar[l]^h \ar[d]^e \\
S & S' \ar[l] & \Spec(K) \ar[l]_g
}
$$
where $X' = X \times_S S'$, $Y = X' \times_{S'} \Spec(K)$,
$K$ is a field, and $\mathcal{F}$ is an abelian sheaf
on $\Spec(K)$ annihilated by $n$ such that
$(f')^{-1}R^qg_*\mathcal{F} \to R^qh_*e^{-1}\mathcal{F}$
is not an isomorphism.
\end{lemma}

\begin{proof}
Choose a commutative diagram
$$
\xymatrix{
X \ar[d]_f & X' \ar[l] \ar[d]_{f'} & Y \ar[l]^h \ar[d]^e \\
S & S' \ar[l] & T \ar[l]_g
}
$$
with $X' = X \times_S S'$ and $Y = X' \times_{S'} T$ and
$g$ quasi-compact and quasi-separated, and an abelian sheaf
$\mathcal{F}$ on $T_\etale$ annihilated by $n$ such that
the base change map
$(f')^{-1}R^qg_*\mathcal{F} \to R^qh_*e^{-1}\mathcal{F}$
is not an isomorphism. Of course we may and do replace $S'$
by an affine open of $S'$; this implies that $T$ is quasi-compact
and quasi-separated. By Lemma \ref{lemma-base-change-q-injective} we see
$(f')^{-1}R^qg_*\mathcal{F} \to R^qh_*e^{-1}\mathcal{F}$
is injective. Pick a geometric point $\overline{x}$ of $X'$
and an element $\xi$ of $(R^qh_*q^{-1}\mathcal{F})_{\overline{x}}$
which is not in the image of the map
$((f')^{-1}R^qg_*\mathcal{F})_{\overline{x}} \to
(R^qh_*e^{-1}\mathcal{F})_{\overline{x}}$.

\medskip\noindent
Consider a morphism $\pi : T' \to T$ with $T'$ quasi-compact
and quasi-separated and denote $\mathcal{F}' = \pi^{-1}\mathcal{F}$.
Denote $\pi' : Y' = Y \times_T T' \to Y$ the base change of $\pi$
and $e' : Y' \to T'$ the base change of $e$. Picture
$$
\vcenter{
\xymatrix{
X' \ar[d]_{f'} & Y \ar[l]^h \ar[d]^e & Y' \ar[l]^{\pi'} \ar[d]^{e'} \\
S' & T \ar[l]_g & T' \ar[l]_\pi
}
}
\quad\text{and}\quad
\vcenter{
\xymatrix{
X' \ar[d]_{f'} & & Y' \ar[ll]^{h' = h \circ \pi'} \ar[d]^{e'} \\
S' & & T' \ar[ll]_{g' = g \circ \pi}
}
}
$$
Using pullback maps we obtain a canonical commutative diagram
$$
\xymatrix{
(f')^{-1}R^qg_*\mathcal{F} \ar[r] \ar[d] &
(f')^{-1}R^qg'_*\mathcal{F}' \ar[d] \\
R^qh_*e^{-1}\mathcal{F} \ar[r] &
R^qh'_*(e')^{-1}\mathcal{F}'
}
$$
of abelian sheaves on $X'$. Let $P(T')$ be the property
\begin{itemize}
\item The image $\xi'$ of $\xi$ in
$(Rh'_*(e')^{-1}\mathcal{F}')_{\overline{x}}$ is
not in the image of the map
$(f^{-1}R^qg'_*\mathcal{F}')_{\overline{x}} \to
(R^qh'_*(e')^{-1}\mathcal{F}')_{\overline{x}}$.
\end{itemize}
We claim that hypotheses (1), (2), and (3) of
Lemma \ref{lemma-formal-argument} hold for $P$
which proves our lemma.

\medskip\noindent
Condition (1) of Lemma \ref{lemma-formal-argument}
holds for $P$ because the \'etale topology of a scheme
and a thickening of the scheme is the same. See
Proposition \ref{proposition-topological-invariance}.

\medskip\noindent
Suppose that $I$ is a directed set and that $T_i$
is an inverse system over $I$ of quasi-compact and quasi-separated
schemes over $T$ with affine transition morphisms.
Set $T' = \lim T_i$. Denote $\mathcal{F}'$ and $\mathcal{F}_i$
the pullback of $\mathcal{F}$ to $T'$, resp.\ $T_i$. Consider
the diagrams
$$
\vcenter{
\xymatrix{
X \ar[d]_{f'} & Y \ar[l]^h \ar[d]^e & Y_i \ar[l]^{\pi_i'} \ar[d]^{e_i} \\
S & T \ar[l]_g & T_i \ar[l]_{\pi_i}
}
}
\quad\text{and}\quad
\vcenter{
\xymatrix{
X \ar[d]_{f'} & & Y_i \ar[ll]^{h_i = h \circ \pi_i'} \ar[d]^{e_i} \\
S & & T_i \ar[ll]_{g_i = g \circ \pi_i}
}
}
$$
as in the previous paragraph. It is clear that $\mathcal{F}'$ on
$T'$ is the colimit of the pullbacks of $\mathcal{F}_i$ to $T'$
and that $(e')^{-1}\mathcal{F}'$ is the colimit of the pullbacks
of $e_i^{-1}\mathcal{F}_i$ to $Y'$.
By Lemma \ref{lemma-relative-colimit-general}
we have
$$
R^qh'_*(e')^{-1}\mathcal{F}' = \colim R^qh_{i, *}e_i^{-1}\mathcal{F}_i
\quad\text{and}\quad
(f')^{-1}R^qg'_*\mathcal{F}' = \colim (f')^{-1}R^qg_{i, *}\mathcal{F}_i
$$
It follows that if $P(T_i)$ is true for all $i$, then
$P(T')$ holds. Thus condition (2) of Lemma \ref{lemma-formal-argument}
holds for $P$.

\medskip\noindent
The most interesting is condition (3) of Lemma \ref{lemma-formal-argument}.
Assume $T'$ is a quasi-compact and quasi-separated scheme over $T$
such that $P(T')$ is true.
Let $Z \subset T'$ be a closed subscheme with complement $V \subset T'$
quasi-compact. Consider the diagram
$$
\xymatrix{
Y' \times_{T'} Z \ar[d]_{e_Z} \ar[r]_{i'} &
Y' \ar[d]_{e'} &
Y' \times_{T'} V \ar[l]^{j'} \ar[d]^{e_V} \\
Z \ar[r]^i &
T' &
V \ar[l]_j
}
$$
Choose an injective map $j^{-1}\mathcal{F}' \to \mathcal{J}$
where $\mathcal{J}$ is an injective sheaf of $\mathbf{Z}/n\mathbf{Z}$-modules
on $V$. Looking at stalks we see that the map
$$
\mathcal{F}' \to \mathcal{G} = j_*\mathcal{J} \oplus i_*i^{-1}\mathcal{F}'
$$
is injective. Thus $\xi'$ maps to a nonzero element of
\begin{align*}
& \Coker\left(
((f')^{-1}R^qg'_*\mathcal{G})_{\overline{x}}
\to
(R^qh'_*(e')^{-1}\mathcal{G})_{\overline{x}}
\right) = \\
&
\Coker\left(
((f')^{-1}R^qg'_*j_*\mathcal{J})_{\overline{x}}
\to
(R^qh'_*(e')^{-1}j_*\mathcal{J})_{\overline{x}}
\right) \oplus \\
& \Coker\left(
((f')^{-1}R^qg'_*i_*i^{-1}\mathcal{F}')_{\overline{x}}
\to
(R^qh'_*(e')^{-1}i_*i^{-1}\mathcal{F}')_{\overline{x}}
\right)
\end{align*}
by part (2) of Lemma \ref{lemma-base-change-q-injective}.
If $\xi'$ does not map to zero in the second summand, then
we use
$$
(f')^{-1}R^qg'_*i_*i^{-1}\mathcal{F}' =
(f')^{-1}R^q(g' \circ i)_*i^{-1}\mathcal{F}'
$$
(because $Ri_* = i_*$ by
Proposition \ref{proposition-finite-higher-direct-image-zero}) and
$$
R^qh'_*(e')^{-1}i_*i^{-1}\mathcal{F} =
R^qh'_*i'_*e_Z^{-1}i^{-1}\mathcal{F} =
R^q(h' \circ i')_*e_Z^{-1}i^{-1}\mathcal{F}'
$$
(first equality by
Lemma \ref{lemma-finite-pushforward-commutes-with-base-change}
and the second because
$Ri'_* = i'_*$ by
Proposition \ref{proposition-finite-higher-direct-image-zero})
to we see that we have $P(Z)$.
Finally, suppose $\xi'$ does not map to zero in the first summand.
We have
$$
(e')^{-1}j_*\mathcal{J} = j'_*e_V^{-1}\mathcal{J}
\quad\text{and}\quad
R^aj'_*e_V^{-1}\mathcal{J} = 0, \quad a = 1, \ldots, q - 1
$$
by $BC(f, n, q - 1)$ applied to the diagram
$$
\xymatrix{
X \ar[d]_f & Y' \ar[l] \ar[d]_{e'} & Y \ar[l]^{j'} \ar[d]^{e_V} \\
S & T' \ar[l] & V \ar[l]_j
}
$$
and the fact that $\mathcal{J}$ is injective.
By the relative Leray spectral sequence for $h' \circ j'$
(Cohomology on Sites, Lemma \ref{sites-cohomology-lemma-relative-Leray})
we deduce that
$$
R^qh'_*(e')^{-1}j_*\mathcal{J} =
R^qh'_*j'_*e_V^{-1}\mathcal{J}
\longrightarrow
R^q(h' \circ j')_* e_V^{-1}\mathcal{J}
$$
is injective. Thus $\xi$ maps to a nonzero element of
$(R^q(h' \circ j')_* e_V^{-1}\mathcal{J})_{\overline{x}}$.
Applying part (3) of Lemma \ref{lemma-base-change-q-injective}
to the injection $j^{-1}\mathcal{F}' \to \mathcal{J}$
we conclude that $P(V)$ holds.
\end{proof}

\end{document}
